\documentclass[10pt,a4paper]{amsart}
\usepackage[margin=19mm]{geometry}
\usepackage{amsmath,amssymb,mathtools}
\usepackage[scr=rsfs,cal=euler]{mathalfa}
\usepackage{xfrac}
\usepackage[unicode,hidelinks,bookmarks=false]{hyperref}
\newcommand{\ie}{i.e.\ }
\newcommand{\cf}{cf.\ }
\newcommand{\resp}{respectively\ }
\newcommand{\Subsection}{\subsection}
\providecommand{\nobreakdash}{\nobreak\hbox{-}}
\setlength{\emergencystretch}{2em}
\multlinegap0pt
\usepackage{bm}
\usepackage{bbm}
\usepackage{dsfont}
\usepackage{xspace}
\usepackage{cancel}
\usepackage{mathtools}
\usepackage{amsthm}
\makeatletter
\@ifundefined{theorem}{\newtheorem{theorem}{Theorem}}{}
\makeatother
\title[Know Your Limits: Entropy Estimation Modeling for Compressio]{Know Your Limits: Entropy Estimation Modeling for Compression and Generalization}
\author{Benjamin L. Badger, Matthew Neligeorge}
\date{Source: arXiv:2511.10618v1 (CC BY 4.0)}
\begin{document}
\maketitle

\noindent\emph{Source and licence.} Know Your Limits: Entropy Estimation Modeling for Compression and Generalization. Version arXiv:2511.10618v1 (CC BY 4.0), distributed under the Creative Commons Attribution 4.0 International licence, as recorded in the arXiv licence metadata for this exact version and shown on the abstract page https://arxiv.org/abs/2511.10618v1. Full rights notice: licenses/arxiv-2511.10618-CC-BY-4.0.txt.\par\smallskip

\noindent\textit{Editorial scope.} Counted unit: the Theorem whose source label is \texttt{None} in the article (source0.tex lines 398--400, source label \texttt{None}), delivered together with the article's own complete original proof of it (source0.tex lines 402--429, one \texttt{proof} environment of 28 lines). No mathematics is written, repaired or re-derived: statement and proof are the article's own text line for line. Required context retained uncounted: none. Every definition, display and label of those lines is retained unchanged. Omitted from this package and not redistributed: 1--397, 401--401, 430--658. Nothing inside a retained excerpt is omitted or altered: every retained body is delivered exactly as the article's lines are written, so no omission receipt of any kind accompanies this package. Citations: the retained excerpts carry no citation command at all, so no bibliography entry of the article is transcribed and no reference is redistributed. Reference adaptation: none was needed and none was made; every reference inside the delivered unit resolves inside it, so no unresolved reference or citation is produced. Printed slips kept literal: no printed word, symbol, hypothesis or number of the counted statement or of its proof was altered. Layout and class: the article's own class is \texttt{article} with packages including geometry, lipsum, enumitem, inputenc, fontenc, url, booktabs, amsfonts, nicefrac, microtype, graphicx, natbib, doi, float; the wrapper is the lane's pinned \texttt{amsart} editorial wrapper at 10pt on A4, which restates the theorem-like environments the retained text needs and adds the article's own macro definitions that the retained text needs (none), with extra wrapper packages (bm, bbm, dsfont, xspace, cancel, mathtools). The delivered numbering is the unit's own. Assets: no retained span contains a figure environment, a graphics-inclusion command, a PDF-page inclusion, a table or a TikZ picture; no image file is copied into this package, the package declares no asset dependency and no image file is redistributed with it. Licence: version arXiv:2511.10618v1 of this article is distributed under the Creative Commons Attribution 4.0 International licence, as recorded in the arXiv licence metadata for this exact version and shown on the abstract page https://arxiv.org/abs/2511.10618v1. A full rights notice is delivered at licenses/arxiv-2511.10618-CC-BY-4.0.txt. Characters: every retained excerpt is pure ASCII, so the delivered engine, XeLaTeX with its default Latin Modern text font, reports no missing character.
    \begin{theorem}
        Generalization decreases when models are trained to minimize their losses below the entropy of the training dataset.
    \end{theorem}

    \begin{proof}

    We assume that some random variable $X$ exists and contains independent and identically distributed disjoint subsets, $X_{train}$ and $X_{text}$, where $\{x' \sim X_{train}\} \; \bigcup \; \{x'' \sim X_{test}\} = \{x \sim X\}$ and $\{x' \sim X_{train}\} \; \bigcap \; \{x'' \sim X_{test}\} = \emptyset$. For notational clarity, we substitute $S = X_{train}$ and $T = X_{test}$ and use $H_B(A)$ to denote the cross-entropy between random variables $A$ and $B$, and $H(A, B)$ to denote the joint entropy of $A$ and $B$. 
    
    We assume that the entropy of the random variable is nonzero, $H(X) > 0$, and by the independence of $S, T$ we have $H(X) = H(S, T) = H(S) + H(T)$.  We also assume that model losses for samples drawn from these random variables, $s \sim S, t \sim T$ are cross-entropies on causal predictions, where we denote the cross-entropy of a model's output on an arbitrarily large number of samples of a random variable as $H_S(O(S, \theta)) = \sum_n \sum_i\Bbb L(O(s_{[:i-1]}, \theta), s_i)_n = \sum_n H_s(O(s_n, \theta), s_n) :  s_n \sim S$ and the same for $T$.

    From Gibbs' inequality, we know that $H_X(O(X, \theta)) \geq H(X) \; \forall  X, \theta$. We decompose $X$ according to the independence of its subsets $S, T$ as follows:
    
    \begin{equation}
        \begin{split}
         H_X(O(X, \theta)) = & H_S(O(S, \theta)) + H_T(O(T, \theta)) \geq H(X) = H(S) + H(T) \implies \\
         & H_S(O(S, \theta)) - H(S) + \left( H_T(O(T, \theta)) - H(T)  \right) \geq 0 
         \end{split}
         \label{eq15}
    \end{equation}

    Therefore as $\theta$ is modified such that $H_S(O(S, \theta)) - H(S)$ decreases below $0$, $H_T(O(T, \theta)) - H(T)$ increases so that the Inequality \ref{eq15} holds. We assumed that $H(T)$ is fixed, so necessarily $H_T(O(T, \theta))$ increases as $H_S(O(S, \theta))$ decreases below $H(S)$. 
    
    In fact, due to the identical distribution of $S, T$ we know that $H(S) = H(T)$ and can put a lower bound on the decrease of generalization as shown in Equation \ref{eq16}.

    \begin{equation}
        H_S(O(S, \theta) - H(S) \geq - \left( H_T(O(T, \theta)) - H(S) \right)
        \label{eq16}
    \end{equation}
    
    We conclude by restating Equation \ref{eq16} in words: as it was assumed $H_S(O(S, \theta)) < H(S)$ making the left side of the inequality negative, the difference between the model's cross-entropy loss and the training dataset's entropy is at least as large as the difference between the model's loss on the test set and that entropy.
    
    \end{proof}

\end{document}
